\section{The actual uniform integral theta denominator}
\label{sec:actual-theta-denominator}

The integral theta estimate now applies to the constructed arithmetic
cover.  Its value identity is already part of the original raw-theta
contract: it is a finite torsion sum, with all its indices and scalars.
The passage from test rings to the actual valuation tests introduces
no analytic comparison or extra denominator.

Use any original $A=1$ deep cover $Y$ satisfying the source hypotheses,
its coefficient field $c.F$, higher half-level $M$, and integral cusp
lattice.  For the actual construction, existence of this cover and
its exact complex normalization is supplied by
Corollary~\ref{cor:actual-good-sized-cover}.  Allow every connection
\texttt{nab} satisfying the original \texttt{IsNab}; its harmonic
values agree with the constructed canonical chart values.  Put
\[
 f=[E:\mathbb Q]/2,\quad r_j=p+q,\quad
 d_a=\texttt{DF}(\mathrm{Dw},a),\quad
 \Delta=|\operatorname{disc}(E)|,\quad
 \kappa=\texttt{sizeK}(\mathrm{Dw},a,n).
\]
Here $n$ is the original even theta degree and $a$ the original
nonnegative harmonic row weight.  Retain the original rational space,
fractional lattice $\texttt{AA}L_0$, tangent eigenspace lattices,
primitive line and all source conditions \texttt{CovGmin.Hyp}.
For the numerical conclusion we use its bounds
\[
 L_b\ge1,\quad \log M\le30L_b,\quad
 \operatorname{fieldSz}(E)\le3L_b,\quad f,R\ge1,
\]
where $R=\texttt{D.R.r}$ is the original size rank.  All other
conditions in that contract, including the weight range and sized
cover, remain unchanged.

\begin{theorem}[The original integral raw-theta supplier]
\label{thm:actual-theta-denominator}
One integer, independent of the individual theta generator, suffices:
\begin{equation}\label{eq:actual-theta-clearing-integer}
 N_A=Z= M^{f r_j n+d_a}\Delta^{d_a}
       =\texttt{intZ}(D,\mathrm{Dw},l_2,n,a)>0,
 \qquad
 \log Z\le67(fRL_b)^2\kappa.
\end{equation}
More precisely, let $\Lambda\subset E^n$ be the original unimodular
lattice, trivial at $\texttt{badC}(Y)$ and at $M$, and let the original
odd $N$ satisfy
\[
 M\mid N,\qquad N\mathcal O_E^n\subseteq
 \Lambda_0:=\Lambda\cap\mathcal O_E^n.
\]
Let $P\in\texttt{harmL}(c,\Lambda,a)$ and let
$\theta\in Y.V(\mathrm{wtZ}(a,n))$ have the original
\texttt{thLaValAt} value on every original test-ring framed point,
with level indicator $x_0\in(\mathrm{LevQc})^n$ and its prescribed
finite Fourier transform.  For every $s\in c.F$ with
\begin{equation}\label{eq:actual-theta-determinant-scalar}
 (s)\prod_{\tau:E\hookrightarrow\mathbb C}
          I_\tau(\Lambda_0)^{m_\tau}
       \subseteq
       \prod_{\tau:E\hookrightarrow\mathbb C}
          I_\tau(\Lambda)^{m_\tau},
 \qquad m_\tau=\mathrm{Dw.mult}(\tau),
\end{equation}
one has
\[
                 Zs\theta\in\texttt{AllWt.latW}
                          (Y.U,\mathrm{wtZ}(a,n)).
\]
This is the complete original
\texttt{CovInt.integral\_smul\_thLa} implication, with $Z$ in place
of its permitted $Z^2$.  No restriction on the valuations of $M$,
on the splitting of residue primes in $E$, or on the ramification of
the coefficient field is added.
\end{theorem}

\begin{proof}
\emph{The actual torsion expression at every valuation test.}
The definition \texttt{thLaValAt} is literally
\texttt{ShortValAt} with scalar parameter $\beta=1$
\cite{actual-denominator-source}.  Write
\[
 \mathcal Q=(\mathrm{LevQ})^n,\qquad C=|\mathcal Q|=M^{f r_j n},
 \qquad I=[\Lambda_0:N\mathcal O_E^n]^{r_j}.
\]
The torsion group $\mathcal Z$ has the full cardinality
$[\Lambda:N\mathcal O_E^n]^{r_j}$; it is not replaced by a
subgroup of cardinality $I$.

Take any of the original finite valuation tests, represented by an
integral framed point $z$ over a complete discrete valuation ring.
The actual \texttt{AllTest} supplier in
Theorem~\ref{thm:actual-level-cover} presents its generic framed
point as a framed base change of an original reduced finite-type
test ring.  Apply the assumed raw-theta identity there.  Faithful
flatness of its witness cover is preserved by base change.  The
canonical line comparisons, torsion equations and entire
annihilator-group enumeration pull back, and the canonical
connection's harmonic values commute with that base change.
The value of $\theta$ pulls back as well, since it is the same
algebraic section on the framed-point functor.

This gives actual finite torsion data over an extension of the
generic field.  They may be taken over a finite algebraic extension:
all the points are killed by $N^2$, so their coordinates lie in
finite \'etale generic-fibre torsion schemes.  The sections are
the unique prescribed normalized pullback and two ordered symmetric
theta translations of the given theta generator; they are defined
over the same finite extension.  If the faithful-flat witness
initially has a larger residue field, retain only this finite field
of definition.  Uniqueness of these section constructions and
functoriality of harmonic evaluation preserve the equality.

Thus, with $d=\deg P\le d_a$, the actual identity is exactly
\begin{equation}\label{eq:actual-theta-torsion-value}
 \theta(z)=\frac1{CIN^d}
    \sum_{q\in\mathcal Q}\zeta_M^{\langle x_0,q\rangle}
    \sum_{u\in\mathcal Z}
         H\!\left(P;(U_N(u_b,t_{q,b};\theta_z))_b\right).
\end{equation}
Here $Nt_{q,b}$ is the original half-level point, $t_{q,b}$ is
killed by a power of the original half-level ideal, and
$U_N=T_{t_{q,b}}^{(N^2)}T_{u_b}^{(N^2)}[N]^\natural\theta_z$
has the original order and pullback normalization.  The Fourier
coefficient is the prescribed root of unity: summing the class
indicator in \texttt{phiHat} leaves exactly this one term.
If one instead writes the original short rule
$g=\beta_H^{-1}\theta$, its defining identity has the extra factor
$\beta_H^{-1}$ on the right.  Multiplication by the same
$\beta_H$ gives \eqref{eq:actual-theta-torsion-value}; no
determinant-ideal factor has been discarded.

\emph{Integral models and the two prime branches.}
The actual model and primitive-line supplier, followed by
Theorem~\ref{thm:actual-good3}, gives the required integral chart
after an unramified extension.  The original separate tangent
eigenlattices remain unchanged.  Reached-section regularity and
the discriminant theorem apply to these actual sections and give
the original loss $\rho\le d_a\nu(\Delta)$.

At a rational residue prime $\ell\mid M$, local triviality of
$\Lambda$ makes both determinant ideals in
\eqref{eq:actual-theta-determinant-scalar} units, so $\nu(s)\ge0$.
Theorem~\ref{thm:local-level-theta-denominator} bounds each full
inner sum after multiplication by $M^d I^{-1}N^{-d}$ below by
$-\rho$, with every exponent $v_\ell(M)$ retained.  Its primary
group cancels the exact local part of $I$ by invariance of the
\emph{entire translated product section} before evaluation.

At $\ell\nmid M$, each $t_{q,b}$ has order prime to $\ell$.
Theorem~\ref{thm:full-primary-theta-trace}, applied after the
prime-to-$\ell$ translates, bounds
$sI^{-1}N^{-d}$ times the same full inner sum below by $-\rho$.
The adapted lattice determinant comparison uses precisely
\eqref{eq:actual-theta-determinant-scalar}; it does not replace
$s$ by an integral scalar.  Primary and prime-to-$\ell$ lifts
commute with their original scalar one.  These are exactly the
two branches of the all-place argument in
Section~\ref{sec:geometry-mixed-level}, whose model and chart
premises have now been supplied.

Multiplication by $\Delta^{d_a}M^{d_a}$ clears both bounds, and
the factor $C$ cancels the entire finite Fourier average in
\eqref{eq:actual-theta-torsion-value}.  All Fourier coefficients
are integral roots of unity.  Under an intervening finite
extension of ramification index $e'$, the raw valuation and its
normalizing index both multiply by $e'$, so the original
normalized valuation is unchanged.  The conclusion therefore
holds in every original valuation test.  Intersecting those
conditions is exactly membership in the original cusp lattice.

\emph{The already prescribed uniform size.}
The exact weight identity and the CM-type cardinality give
\[
 \kappa=d_a+f r_j n/2,
 \qquad f r_j n+d_a\le2\kappa,
 \qquad d_a\le\kappa.
\]
The definition of field size gives
$\log\Delta\le[E:\mathbb Q]\operatorname{fieldSz}(E)
\le6fL_b\le7fL_b$, retaining the allowance in the original
denominator estimate.
Consequently
\[
 \log Z
  =(f r_j n+d_a)\log M+d_a\log\Delta
  \le60L_b\kappa+7fL_b\kappa
  \le67fL_b\kappa
  \le67(fRL_b)^2\kappa.
\]
All parameters are the original ones.  The integer is fixed before
$x_0,\Lambda,N,P,s$ and the individual theta section are chosen.
\end{proof}

The remaining short-section inputs are the actual theta spanning
family and its archimedean norm estimate.  The theorem supplies their
finite-place lattice bound; it does not infer either of those other
inputs from integrality.
